% Canonical core-geometry figure. Requires tikz-cd.
% Lean correspondence: HCP.CoupledGeometry; HCP.CoupledSemantics.
\begin{figure}[htbp]
\centering
\begin{tikzcd}[row sep=2.8em, column sep=3.2em]
 & \mathcal D_{\mathbb R}
     \arrow[dl, "M_B"'] \arrow[dr, "M_C"] & \\
 C^1(F_B) \arrow[rr, "L_{B\to C}"] \arrow[d, "P_B"']
   && C^1(F_C) \arrow[d, "P_C"] \\
 W_B \arrow[rr, "S_{B,C}|_{W_B}"]
     \arrow[dr, "\operatorname{info}_B"']
   && W_C \arrow[dl, "\operatorname{info}_C"] \\
 & \mathbb R^n &
\end{tikzcd}

\medskip
\begin{tikzcd}[row sep=2.2em, column sep=3.2em]
 \mathcal I(K_B) \arrow[rr, "\simeq_{\mathrm H}"]
       \arrow[dr, "\operatorname{Th}"']
   && \mathcal I(K_C) \arrow[dl, "\operatorname{Th}"] \\
 & \alpha &
\end{tikzcd}
\caption{The core recovery diagram and its semantic class.
Here \(\mathcal D_{\mathbb R}=\mathbb R\oplus\mathbb R^n\oplus\mathbb R^m\),
\(W_B=\ker\Delta_B\), and \(\operatorname{info}_B(0,z,Bz)=z\).
The upper triangle assembles the same coordinates, the square projects to
harmonic messages, and the lower triangle reads their coefficients.
The separate semantic triangle uses the actual cochain complexes
\(K_B=C^\bullet(F_B)\), their typed interpretations, and the common declared
homological theory \(\alpha\). The semantic equivalence is induced by
\((1,S_{B,C},1)\). Its shear acts on harmonic messages in the
large diagram; native transport acts on ciphertexts.}
\label{fig:coupled-core-diagrams}
\end{figure}
